\documentclass[11pt]{article}
\usepackage{fontspec}
\usepackage[french]{babel}
\usepackage{amsmath}
\usepackage{amssymb}
\usepackage{amsthm}
\usepackage{mathrsfs}
\usepackage{tikz-cd}
\usepackage{enumitem}

\newcommand{\printedpage}[1]{\marginpar{\footnotesize p.~#1}}
\newcommand{\msp}[1]{\marginpar{\footnotesize [#1]}}
% underlined roman name of a category (Fibhot, Cofib, Parf, ...)
\newcommand{\usl}[1]{\underline{\mathrm{#1}}}
% encircled roman numeral of the manuscript
\newcommand{\encadre}[1]{\tikz[baseline=(n.base)]\node[draw,rounded corners=3pt,inner sep=1.5pt](n){#1};}

\begin{document}

\printedpage{49}
% head: A. Grothendieck : Dérivateurs, Ch. VI, Édition des Universités de Montpellier II et Paris VII
% foot: Version 9 novembre 2001

\[
\begin{tikzcd}
\mathcal{F}'\Sigma^{-1} \arrow[r, "\approx"] & \mathcal{F}\Sigma^{-1} \arrow[l, bend left=40, "\approx"', "\Psi_{\mathcal{F}',\mathcal{F}}"]
\end{tikzcd}
\qquad\qquad (\Sigma = W_S \text{ ou } W_S^{\mathrm{g}})
\]

pour deux sous-catégories pleines $\mathcal{F}'$, $\mathcal{F}$
\[
\usl{Cofib}\,S \subseteq \mathcal{F}' \subseteq \mathcal{F} \subseteq \mathrm{Cat}/S.
\]
Les foncteurs $\Phi$, $\Psi$ envoient la sous-catégorie pleine $\usl{Fibhot}\,S$ de ‘l’étage’, dans les sous-catégories $\usl{Fibsc}_0S$, $\usl{Cofibsc}_0S$, respectivement, donc induisent à l’étage des foncteurs $\Phi_0$, $\Psi_0$
\[
\usl{Fibsc}_0S \xleftarrow{\Phi_0} \usl{Fibhot}\,S \xrightarrow{\Psi_0} \usl{Cofibsc}_0S,
\]
jouant à l’étage exactement le même rôle que $\Phi$, $\Psi$ au rez-de-chaussée. La seule différence, c’est que à l’étage on a $W_S = W_S^{\mathrm{d}} = W_S^{\mathrm{g}}$ $(= W_S^{\mathrm{f}})$,
\msp{page 78}
donc il n’y a qu’une suite de catégories localisées à examiner. On trouve ainsi
\[
\begin{tikzcd}
(\usl{Fibsc}_0S)W_S^{-1} \arrow[r, "\approx"] & \mathcal{F}W_S^{-1} \arrow[l, bend left=40, "\approx"', "\Phi_{0,\mathcal{F}}"]
\end{tikzcd}
\]
\[
\begin{tikzcd}
\mathcal{F}'W_S^{-1} \arrow[r, "\approx"] & \mathcal{F}W_S^{-1} \arrow[l, bend left=40, "\approx"', "\Phi_{0,\mathcal{F}',\mathcal{F}}"]
\end{tikzcd}
\]
pour $\mathcal{F}$ (ou $\mathcal{F}$, $\mathcal{F}'$) sous-catégorie(s) pleine(s) de $\usl{Fibhot}\,S$ telle(s) que
\[
\usl{Fib}_0S \subseteq \mathcal{F} \subseteq \usl{Fibhot}\,S \qquad\text{resp.}\qquad \usl{Fib}_0S \subseteq \mathcal{F}' \subseteq \mathcal{F} \subseteq \usl{Fibhot}\,S,
\]
et de même du côté droit pour les sous-catégories pleines $\mathcal{F}$ ou $\mathcal{F}$, $\mathcal{F}'$ de $\usl{Fibhot}\,S$ satisfaisant
\[
\usl{Cofib}_0S \subseteq \mathcal{F} \subseteq \usl{Fibhot}\,S \qquad\text{resp.}\qquad \usl{Cofib}_0S \subseteq \mathcal{F}' \subseteq \mathcal{F} \subseteq \usl{Fibhot}\,S,
\]
NB. La catégorie $\usl{Parf}\,S$ contient à la fois $\usl{Fib}_0S$ et $\usl{Cofib}_0S$, et fait donc partie des catégories admissibles tout comme $\mathcal{F}$, ou [?] $\mathcal{F}'$.

Ainsi, les quatre foncteurs $\Phi$, $\Psi$, $\Phi_0$, $\Psi_0$ donnent tous les quasi-inverses qu’on peut souhaiter au rez-de-chaussée d’une part ($\Phi$ et $\Psi$), à l’étage de l’autre ($\Phi_0$, $\Psi_0$). Ils ramènent donc toutes les localisées pertinentes aux quatre suivantes
\[
\begin{gathered}
(\mathrm{Cat}/S)W_S^{-1},\ \ (\mathrm{Cat}/S)(W_S^{\mathrm{d}})^{-1},\ \ (\mathrm{Cat}/S)(W_S^{\mathrm{g}})^{-1} \qquad \text{(au rez-de-chaussée)}\\
(\usl{Fibhot}\,S)W_S^{-1} \qquad \text{(à l’étage).}
\end{gathered}
\]

\encadre{VII} Il reste à comparer la localisée à l’étage avec les localisées au rez-de-chaussée. Les flèches verticales descendantes de l’étage au rez-de-chaussée, sont compatibles avec $W_S^{\mathrm{d}}$ du côté gauche du diagramme, avec $W_S^{\mathrm{g}}$ du côté droit, donc elles induisent essentiellement \textit{deux} foncteurs différents
\printedpage{50}
% head: A. Grothendieck : Dérivateurs, Ch. VI, Édition des Universités de Montpellier II et Paris VII
% foot: Version 9 novembre 2001
\msp{page 79}
$i$, $j$ sur $(\usl{Fibhot}\,S)W_S^{-1}$, dans $(\mathrm{Cat}/S)(W_S^{\mathrm{d}})^{-1}$ et $(\mathrm{Cat}/S)(W_S^{\mathrm{g}})^{-1}$ :
\[
\begin{tikzcd}[column sep=small, row sep=large]
& (\usl{Fibhot}\,S)W_S^{-1} \arrow[dl, "i"'] \arrow[dr, "j"] \arrow[dd, dotted, "\approx"', "\delta"] & \\
\mathrm{HOT}_W(S) \overset{\text{déf}}{\approx} (\mathrm{Cat}/S)(W_S^{\mathrm{d}})^{-1} \arrow[dr, dotted, "\rho"'] & & (\mathrm{Cat}/S)(W_S^{\mathrm{g}})^{-1} \overset{\text{déf}}{\approx} \mathrm{HOT}^{o}_{W}(S) \arrow[dl, dotted, "\sigma"] \\
& \begin{array}{l} (\mathrm{Cat}/S)W_S^{-1} \\ \overset{\text{déf}}{\approx} \mathrm{HOT}^{\mathrm{lc}}_{W}(S) \end{array} &
\end{tikzcd}
\qquad ,
\]
de façon à faire commuter le carré complété par les pointillés, représentant les foncteurs de localisation $\rho$, $\sigma$. On a bien sûr
\[
\begin{aligned}
(*) \qquad \rho i = \sigma j : (\usl{Fibhot}\,S)W_S^{-1} \longrightarrow (\mathrm{Cat}/S)W_S^{-1}\\
\text{induit par l’inclusion}\\
\delta : \usl{Fibhot}\,S \hookrightarrow \mathrm{Cat}/S
\end{aligned}
\]
Nous avons énoncé dans \encadre{I} que cette flèche est une équivalence, de sorte que $\rho$ et $\sigma$ peuvent être interprétées, à isomorphisme de foncteurs près, comme des rétractions de $\mathrm{HOT}_W(S)$ et de $\mathrm{HOT}^{o}_{W}(S)$ sur une sous-catégorie commune, s’interprétant soit comme $(\usl{Fibhot}\,S)W_S^{-1}$, soit comme $(\mathrm{Cat}/S)W_S^{-1}$. Il reste à décrire le foncteur quasi-inverse de $(*)$, ce qui donnera en même temps les quasi-inverses de tous les foncteurs verticaux (de l’étage vers le rez-de-chaussée) dans le diagramme p. 76, plus exactement pour leurs localisés, pour $W_S$.

\msp{page 80}
C’est ce qui est obtenu à l’aide du foncteur ‘montant’ (du rez-de-chaussée vers l’étage) $C$, qu’on peut regarder comme un foncteur $\mathrm{Cat}/S \longrightarrow \usl{Fibhot}\,S$, mais qui en fait prend ses valeurs dans $\usl{Parf}\,S$. (C’est là un détail sans grande importance, me semble-t-il.)
\[
\mathrm{Cat}/S \xrightarrow{C} \usl{Parf}\,S \hookrightarrow \usl{Fibhot}\,S.
\]
Ce foncteur est compatible avec les localiseurs $W_S$ des deux côtés, et établit des équivalences (quasi-inverses des équivalences
\[
\begin{gathered}
(\usl{Parf}\,S)W_S^{-1} \xrightarrow{\sim} (\mathrm{Cat}/S)W_S^{-1} \qquad\text{et}\qquad (\usl{Fibhot}\,S)W_S^{-1} \xrightarrow{\sim} (\mathrm{Cat}/S)W_S^{-1}),\\
(\mathrm{Cat}/S)W_S^{-1} \xrightarrow{\sim} (\usl{Parf}\,S)W_S^{-1} \left( \xrightarrow{\sim} (\usl{Fibhot}\,S)W_S^{-1} \right).
\end{gathered}
\]

\printedpage{51}
% head: A. Grothendieck : Dérivateurs, Ch. VI, Édition des Universités de Montpellier II et Paris VII
% foot: Version 9 novembre 2001
Donc, quand les trois catégories HOT pour $S$ sont interprétées en termes de localisées de $\mathrm{Cat}/S$, c’est le foncteur $C$ (plus exactement le composé $\mathrm{Cat}/S \xrightarrow{C} \usl{Parf}\,S \hookrightarrow \mathrm{Cat}/S$) qui fournit à la fois le foncteur d’inclusion (!) $i$, et le foncteur d’inclusion $j$ dans $\mathrm{HOT}_W(S)$ et $\mathrm{HOT}^{o}_{W}(S)$ respectivement. Le foncteur en question
\[
C' : \mathrm{Cat}/S \longrightarrow \mathrm{Cat}/S
\]
satisfait bien
\[
C'(W_S) \subseteq W_S^{\mathrm{d}} \cap W_S^{\mathrm{g}}
\]
et définit les deux foncteurs
\[
\begin{tikzcd}[column sep=large]
& (\mathrm{Cat}/S)(W_S^{\mathrm{d}})^{-1} \approx \mathrm{HOT}_W(S) \\
\begin{array}{l} (\mathrm{Cat}/S)W_S^{-1} \\ \approx \mathrm{HOT}^{\mathrm{lc}}_{W}(S) \end{array} \arrow[ur, "i"] \arrow[dr, "j"'] & \\
& (\mathrm{Cat}/S)(W_S^{\mathrm{g}})^{-1} \approx \mathrm{HOT}^{o}_{W}(S)
\end{tikzcd}
\]
Cela nous permet d’expliciter les formules d’adjonction de \encadre{IV}, quand on y interprète $\xi$, $\eta$
\msp{page 81}
comme provenant de $X$, $Y$ dans $\mathrm{Cat}/S$ :
\[
\boxed{\mathrm{Hom}_{W_S^{\mathrm{d}}}(X,CY) \simeq \mathrm{Hom}_{W_S^{\mathrm{g}}}(X,CY) \simeq \mathrm{Hom}_{W_S}(X,Y),}
\]
où $\mathrm{Hom}_{W_S^{\mathrm{d}}}$, $\mathrm{Hom}_{W_S^{\mathrm{g}}}$ et $\mathrm{Hom}_{W_S}$ désignent respectivement les Hom dans les catégories localisées pour $W_S^{\mathrm{d}}$, $W_S^{\mathrm{g}}$ et $W_S$, i.e. dans $\mathrm{HOT}_W(S)$, $\mathrm{HOT}^{o}_{W}(S)$ et $\mathrm{HOT}^{\mathrm{lc}}_{W}(S)$ respectivement (\footnote[51]{Cette formule suggère que si $X \in \mathrm{Ob}\,\mathrm{Cat}/S$, $Y \in \mathrm{Ob}\,\usl{Fibhot}\,S$, alors
\[
\mathrm{Hom}_{W_S^{\mathrm{d}}}(X,Y) \simeq \mathrm{Hom}_{W_S^{\mathrm{g}}}(X,Y) \simeq \mathrm{Hom}_{W_S}(X,Y),
\]
mais ça résulte de la formule d’adjoinction et de $Y \longrightarrow C(Y)$ dans $W_S^{\mathrm{d}} \cap W_S^{\mathrm{g}}$ pour $Y$ dans $\usl{Fibhot}\,S$.}). % [sic] adjoinction -> adjonction, spelling (elsewhere "adjonction")

\textbf{Commentaire.} Toutes les constructions et les énoncés de \encadre{I} à \encadre{VI}, à l’exception du fait que les foncteurs descendants induisent des équivalences pour les localisés \textit{grossiers}, et des formules d’adjonction de \encadre{IV}, sont de nature quasiment soritale, et n’exigent d’autres conditions sur le localiseur fondamental $W \subseteq \mathrm{Fl}\,\mathrm{Cat}$ que L1, L3 quater, L4 bis, récapitulées comme a) b) c) de la page 45. (On n’a même pas à utiliser d).) Il est vrai qu’en l’absence de la condition L5 (page 43), on ignore si à l’étage les deux localiseurs pertinents, $W_S$ et $W_S^{\mathrm{f}}$ (égal à $W_S^{\mathrm{d}}$ du côté gauche, à $W_S^{\mathrm{g}}$ du côté droit) sont égaux. Donc à priori on trouve à l’étage \textit{deux} catégories localisées essentiellement distinctes, et non pas une seule, savoir $(\usl{Fibhot}\,S)\Sigma^{-1}$ avec $\Sigma = W_S$ ou $\Sigma = W_S^{\mathrm{f}}$ $(= W_S^{\mathrm{d}} = W_S^{\mathrm{g}})$.

\printedpage{52}
% head: A. Grothendieck : Dérivateurs, Ch. VI, Édition des Universités de Montpellier II et Paris VII
% foot: Version 9 novembre 2001
Par contre, le fait que les foncteurs de descente induisent des équivalences entre les localisés grossiers est plus profond, il exige (me semble-t-il) qu’on connaisse un foncteur ‘montant’ $C$, disons de $\mathrm{Cat}/S$ vers $\usl{Fibhot}\,S$, compatible avec $W_S$, et muni de
\[
X \xrightarrow{i_X} C(X),
\]
qui soit dans $W_S$ pour $X$ quelconque. Quand on sait que Cat est une catégorie de Quillen, même sans avoir les propriétés des foncteurs $\usl{Ch}_\infty$ ou $\usl{Ch}$ (\footnote[52]{ni même être sûr de l’existence d’une factorisation \textit{fonctorielle} de $X \xrightarrow{f} S$ en $X \xrightarrow{i_X} C(X) \xrightarrow{\bar{f}} S$ avec $i_X \in W$ et $\bar{f}$ parfait.}), on doit sans doute
\msp{page 82}
pouvoir se débrouiller. Mais la formule d’adjonction me semble rester inattendue et profonde, même en admettant les résultats de \textsc{Quillen} et \textsc{Thomason} (\footnote[53]{Voir cependant le théorème scholie 2, p. 108-113. Ce n’est finalement pas si profond qu’il me semblait, mais quasiment sorital!}). Elle parait liée ici à une propriété assez particulière du foncteur $C$, savoir le lemme p. 68
\[
\begin{gathered}
i_{C(X)} \underset{\big\uparrow}{=} C(i_X) : C(X) \longrightarrow C(C(X)),\\
\scriptstyle \text{dans } \underbrace{(\mathrm{Cat}/S)(W_S^{\mathrm{d}})^{-1}}_{\mathrm{HOT}_W(S)} \text{ et dans } \underbrace{(\mathrm{Cat}/S)(W_S^{\mathrm{g}})^{-1}}_{\mathrm{HOT}^{o}_{W}(S)}
\end{gathered}
\]
ce qui n’a pas l’air purement formel. Il serait intéressant d’examiner dans quelle mesure cette formule garde un sens dans une ‘catégorie de modèles’ $(\mathcal{M},W)$ plus ou moins quelconque, ou provenant d’un triple de Quillen $(W,\mathrm{Cof},\mathrm{Fib})$ dans $\mathcal{M}$. Il est clair que la factorisation $X \xrightarrow{f} S$ en $X \xrightarrow{i_X} C(X) \xrightarrow{\bar{f}} S$ $(i_X \in W,\ \bar{f} \in \mathrm{Parf})$, jouant le rôle des ‘fibrations’ dans le contexte Quillen, a un sens intrinsèque indépendant du choix d’une construction particulière, chose qu’on peut préciser de façon parfaite. C’est la description des localiseurs $W_S^{\mathrm{d}}$, $W_S^{\mathrm{g}}$ dans $\mathcal{M}/S$ (pour $S \in \mathrm{Ob}\,\mathcal{M}$) qui n’est pas comprise de façon générale, en dehors du cas $\mathcal{M} = \mathrm{Cat}$. En d’autre termes, je ne comprends pas comment construire des catégories $\mathrm{HOT}_{\mathcal{M},W}(S)$, $\mathrm{HOT}^{o}_{\mathcal{M},W}(S)$, même si $(\mathcal{M},W)$ provient d’un triple de Quillen, je vois seulement la description d’une catégorie $\mathrm{HOT}^{\mathrm{lc}}_{\mathcal{M},W}(S) = (\mathcal{M}/S)(W_S^{\text{grossière}})^{-1}$, et un foncteur de localisation $\mathcal{M}/S \longrightarrow \mathrm{HOT}^{\mathrm{lc}}_{\mathcal{M},W}(S)$. Il faut arriver à décrire dans $\mathcal{M}/S$ des $W$-équivalences (\textit{‘relatives à $S$’}) plus fines que la grossière donnée par $W$ sur la source $X$ sans plus. Dans Cat j’y arrive en utilisant soit la famille des isomorphismes locaux (pour décrire $W_S^{\mathrm{d}}$), soit celle des isomorphismes colocaux (pour décrire $W_S^{\mathrm{g}}$) - en tous les cas, il y a une notion de ‘localisation’ (soit sur $S$, soit sur $S^{o}$ \dots), avec % [sic] En d’autre termes -> En d’autres termes, fixed expression requires plural
\msp{page 83}
toute la richesse des intuitions liées à une telle notion. Peut-être, sans le savoir ai-je toujours traité Cat non comme une catégorie ‘nue’, mais comme un site, ou plutôt (puisqu’il
\printedpage{53}
% head: A. Grothendieck : Dérivateurs, Ch. VI, Édition des Universités de Montpellier II et Paris VII
% foot: Version 9 novembre 2001
y a l’involution $X \longmapsto X^{o}$ qui ne peut être éludée), comme un ‘bi-site’ - on y a \textit{deux} manières de localiser, soit par les $X/x$ (topologie de Zariski!) soit par les $x\backslash X$ (topologie co-Zariskienne). Il y a aussi la topologie lisse (toute proche sans doute de l’intuition de la cohomologie étale), et la topologie propre - et on s’attend que la lisse donne essentiellement les mêmes résultats que la Zariskienne (vue que la topologie étale lui est égale - il n’y a pas lieu d’en faire un plat \dots); et la topologie propre les mêmes résultats (p. ex. la même cohomologie?) que la topologie co-Zariskienne. % [sic] vue que -> vu que, conjunctional locution

Il semble bien que dans toutes les catégories de modèles utilisées jusqu’à présent pour décrire Hot comme une localisée, on dispose de notions de localisation : (Top), $\mathbf{\Delta}^\wedge$ (ou $A^\wedge$, avec $A$ catégorie-test quelconque), (Cat); et dans Cat on dispose même de \textit{deux} telles notions. Je pense par contre aux $\infty$-groupoïdes, c’est peut-être la seule catégorie de modèles à tel point proche déjà de Hot elle-même, qu’on n’y discerne pas de notion de localisation, et que l’idée même de vouloir s’y localiser (si ce n’est en ‘localisant’ Hot lui-même) paraît quasiment
\msp{page 84}
saugrenue.

Cette rapide digression ne convainc à présent que la compatibilité dite pour la construction de $C(X)$ et $i_X : X \longrightarrow C(X)$, ne peut en aucun cas être purement formelle, une pure histoire de catégories de Quillen disons. Ça donnera au mieux la compatibilité dans $(\mathcal{M}/S)W_S^{-1}$, ce qui est ridicule! Mais si on prend comme localiseur $\Sigma$ dans $\mathcal{M}/S$ les isomorphismes sans plus, la relation de compatibilité va devenir fausse à tous les coups! Comme ensemble de flèches $\Sigma$ un peu plus raisonnable, qui soit encore très petit, on peut songer p. ex. à la $W$-équivalence \textit{universelle} sur $S$ (nettement plus fine encore que $W_S^{\mathrm{d}}$, $W_S^{\mathrm{g}}$ dans le cas particulier de Cat). Aurait-on encore compatibilité pour cet ensemble de flèches? Je suis aussi un peu perplexe par le fait (pas bien compris) que les formules d’adjonction pour les topologies Zariskienne et co-Zariskienne soient en sens opposé. (Ou me serai-je trompé en dualisant à la va-vite?) Mais à vrai dire, cette compatibilité à elle seule, même quand on utilise une structure de site sur $\mathcal{M}$, ne suffit pas à écrire une formule d’adjonction. On a deux catégories localisées, % [sic] ne convainc -> me convainc, the verb needs an object pronoun
\[
\mathrm{HOT}^{\mathrm{lc}}_{\mathcal{M}} = (\mathcal{M}/S)W_S^{-1}, \qquad \mathrm{HOT}_W(S) = (\mathcal{M}/S)(W_S^{\mathrm{f}})^{-1},
\]
où $W_S^{\mathrm{f}}$ désigne une $W$-équivalence ‘fine’, correspondant à la notion de localisation, plus exactement aux changements de base ‘admissibles’ qu’on a choisis, et on a un foncteur de localisation
\msp{page 85}
\[
\mathrm{Hot}_W(S) \longrightarrow \mathrm{Hot}^{\mathrm{lc}}_{W}(S),
\]
mais pas de foncteur en sens inverse. Mais si, mais si, en travaillant avec la catégorie similaire $\usl{Fibhot}\,S$, ça doit garder un sens tout au moins quand on a une bonne notion de ‘fibration’. Je vois un petit travail de dégrossissage conceptuel à faire, avec peut-être une formule d’adjonction à la clef, mais je ne vais pas m’y lancer à présent.

\end{document}
